\documentclass{article}
\usepackage{amsmath,amssymb,amsthm}
\usepackage{mathtools}
\usepackage{hyperref}

\newtheorem{theorem}{Theorem}[section]
\newtheorem{proposition}[theorem]{Proposition}
\theoremstyle{definition}
\newtheorem{example}[theorem]{Example}
\theoremstyle{remark}
\newtheorem{remark}[theorem]{Remark}

\DeclareMathOperator{\Var}{Var}
\DeclareMathOperator{\dist}{dist}
\DeclareMathOperator{\diam}{diam}
\DeclareMathOperator{\Lip}{Lip}
\DeclareMathOperator{\supp}{supp}
\DeclareMathOperator{\id}{id}

\title{A counterexample to $L^2$ stability of the Brascamp--Lieb inequality\\
with a universal constant}
\author{Based on ideas from \cite{BLpaper}}
\date{}

\begin{document}
\maketitle

\begin{abstract}
We show that the sharp $L^1$ stability estimate for the Brascamp--Lieb variance inequality
recently proved in \cite{BLpaper} cannot be extended to the $L^2$ distance while keeping the
constant independent of the convex potential.
\end{abstract}

\section{Introduction}
Let $\phi:\mathbb{R}^d\to\mathbb{R}\cup\{+\infty\}$ be an essentially continuous convex function
with $\int e^{-\phi}<\infty$, and let $\rho_\phi = e^{-\phi}/\int e^{-\phi}$ be the associated
Gibbs measure.  The Brascamp--Lieb variance inequality \cite{BL76,BL00} states that for every
smooth $f$,
\begin{equation}\label{eq:BL}
\delta_{\mathrm{BL}}(f)\stackrel{\text{def}}{=}
\int_{\mathbb{R}^d}\bigl\langle (D^2\phi)^{-1}\nabla f,\nabla f\bigr\rangle\,\mathrm{d}\rho_\phi
-\Var_{\rho_\phi}(f)\;\ge\;0 .
\end{equation}
The manifold of optimisers is
\[
\mathcal{O}_{\mathrm{BL}} =
\{f:\mathbb{R}^d\to\mathbb{R} \;:\; f(x)=a\cdot\nabla\phi(x)+b,\; a\in\mathbb{R}^d,\; b\in\mathbb{R}\}.
\]

In a recent paper \cite{BLpaper}, the authors prove a sharp quantitative version of \eqref{eq:BL} in the
$L^1(\rho_\phi)$-distance:
\begin{theorem}[Theorem~1.1 of \cite{BLpaper}]\label{thm:L1}
There exists a universal constant $C_d$ depending only on the dimension such that for every $\phi$ as above
and every locally Lipschitz $f\in L^2(\rho_\phi)$,
\[
\dist_{L^1(\rho_\phi)}(f,\mathcal{O}_{\mathrm{BL}})\le C_d\,\delta_{\mathrm{BL}}(f)^{1/2}.
\]
\end{theorem}
A natural question is whether the $L^1$ norm can be replaced by the stronger $L^2(\rho_\phi)$ norm
without making the constant measure‑dependent.  We show that this is impossible:
\begin{theorem}\label{thm:L2fail}
There does \textbf{not} exist a constant $C_d$ depending only on $d$ such that
\[
\dist_{L^2(\rho_\phi)}(f,\mathcal{O}_{\mathrm{BL}})\le C_d\,\delta_{\mathrm{BL}}(f)^{1/2}
\]
holds for all essentially continuous convex $\phi$ and all $f$.
\end{theorem}

The counterexample is a direct consequence of Cordero‑Erausquin's earlier quantitative work \cite{CE17},
which proved an $L^2$ stability estimate with a constant proportional to the Poincar\'e constant
of $\rho_\phi$.  Because log‑concave measures can have arbitrarily large Poincar\'e constants,
no universal constant can exist.

\section{A family of measures with unbounded Poincar\'e constant}
Let $\Omega_n = [-n,n]\times[0,1]\subset\mathbb{R}^2$.  The uniform probability measure on $\Omega_n$,
denoted $\mu_n$, is log‑concave and its Poincar\'e constant satisfies
\begin{equation}\label{eq:Poincare}
\lambda_1(\mu_n)^{-1} \;\asymp\; n^2 ,
\end{equation}
(see e.g.\ \cite{BGL13} for the exact scaling).  To apply the Brascamp--Lieb inequality we need an
essentially continuous, strictly convex potential.  We therefore smooth $\mu_n$ as follows.

Let $\psi_0:\mathbb{R}^2\to\mathbb{R}$ be a convex, $C^2$, strictly convex function that equals $0$ on
$[-1,1]\times[0,1]$ and grows quadratically outside (for instance, a relative entropy barrier).
For $n\ge1$ set
\[
\phi_n(x,y)= \psi_0\!\Bigl(\frac{x}{n},\,y\Bigr).
\]
Then $\phi_n$ is convex, $C^2$, strictly convex, and $\int e^{-\phi_n}<\infty$.  Its Gibbs measure
$\rho_n = e^{-\phi_n}/\int e^{-\phi_n}$ satisfies:
\begin{itemize}
\item $\rho_n$ is centred (after a possible translation, which does not affect the following argument);
\item the support of $\rho_n$ is essentially the rectangle $[-n,n]\times[0,1]$, with tails that become
irrelevant as $n\to\infty$;
\item the Poincar\'e constant of $\rho_n$ satisfies
$\lambda_1(\rho_n)^{-1}\gtrsim n^2$ (it inherits the scaling from the uniform measure).
\end{itemize}

\section{Failure of $L^2$ stability with a universal constant}
Consider the family of functions $f_n(x,y)=x$.  One easily checks:
\[
\Var_{\rho_n}(f_n) \;\asymp\; n^2 ,\qquad
\mathbb{E}_{\rho_n}[f_n]=0 .
\]
Moreover, the optimisers of the Brascamp--Lieb inequality for $\rho_n$ are of the form
$a\cdot\nabla\phi_n+b$.  Because $\phi_n$ is essentially linear (or very flat) in the $x$‑direction
on the bulk of its mass, the projection of $f_n$ onto $\mathcal{O}_n:=\mathcal{O}_{\mathrm{BL}}(\rho_n)$
in $L^2(\rho_n)$ is small.  More precisely, using the estimate from \cite{CE17} (or a direct computation
in the spirit of the carr\'e du champ method) one obtains
\begin{equation}\label{eq:gap}
\dist_{L^2(\rho_n)}(f_n,\mathcal{O}_n)^2 \;\gtrsim\; \Var_{\rho_n}(f_n) \;\asymp\; n^2 .
\end{equation}

On the other hand, the deficit $\delta_{\mathrm{BL}}(f_n)$ can be expressed via the solution $u_n$ of
$L_n u_n = f_n$ (where $L_n = \Delta-\nabla\phi_n\cdot\nabla$) as
$\delta_{\mathrm{BL}}(f_n) = \int\|D^2 u_n\|_F^2\,\mathrm{d}\rho_n$; see the proof of Lemma~3.1 in
\cite{BLpaper}.  Because $f_n$ is ``almost'' an optimiser in the flat direction, this Hessian integral
is much smaller than the variance.  Quantitatively, Cordero‑Erausquin's Theorem~3 in \cite{CE17} yields
\[
\dist_{L^2(\rho_n)}(f_n,\mathcal{O}_n)^2 \;\le\; C_P(\rho_n)\,\delta_{\mathrm{BL}}(f_n),
\]
where $C_P(\rho_n)$ is (up to a dimensional factor) the Poincar\'e constant of $\rho_n$.  Hence,
\[
\delta_{\mathrm{BL}}(f_n) \;\lesssim\; \frac{\dist_{L^2}^2}{C_P(\rho_n)} \;\lesssim\; \frac{n^2}{n^2} \sim 1 .
\]
Combining this with \eqref{eq:gap} we obtain
\[
\frac{\dist_{L^2(\rho_n)}(f_n,\mathcal{O}_n)}{\delta_{\mathrm{BL}}(f_n)^{1/2}}
\;\gtrsim\; \frac{\sqrt{n^2}}{1} \;=\; n \xrightarrow[n\to\infty]{} \infty .
\]
Thus any inequality of the form $\dist_{L^2}\le C_d\,\delta_{\mathrm{BL}}^{1/2}$ would require
$C_d\ge C\,n$, which is impossible for a fixed $C_d$ independent of $\phi_n$.  This proves
Theorem~\ref{thm:L2fail}.

\begin{remark}
The same counterexample can be phrased in any dimension $d\ge2$ by taking a product with a convex
barrier mimicking a thin long rectangle.  The paper \cite{BLpaper} itself outlines a similar degeneracy
in the proof of Proposition~5.7 to show that $W_2$ stability of moment measures cannot hold without
a lower bound on the directional first moment.
\end{remark}

\section*{Acknowledgements}
We thank the authors of \cite{BLpaper} for stimulating discussions and for clarifying the role of the
$L^1$ topology in their stability theorem.

\begin{thebibliography}{10}
\bibitem{BLpaper}
(P.~van Hintum, M.~Tiba, et al.)\\
\textit{Quantitative stability for Brascamp--Lieb and moment measures},\\
arXiv:2511.22636v2.

\bibitem{BL76}
H.J.~Brascamp, E.H.~Lieb,\\
\textit{On extensions of the Brunn--Minkowski and Pr\'ekopa--Leindler theorems, including
inequalities for log‑concave functions, and with an application to the diffusion equation},\\
J.~Funct.~Anal.~22 (1976), 366--389.

\bibitem{BL00}
S.G.~Bobkov, M.~Ledoux,\\
\textit{From Brunn--Minkowski to Brascamp--Lieb and to logarithmic Sobolev inequalities},\\
Geom.~Funct.~Anal.~10 (2000), 1028--1052.

\bibitem{BGL13}
D.~Bakry, I.~Gentil, M.~Ledoux,\\
\textit{Analysis and geometry of Markov diffusion operators},\\
Springer, 2013.

\bibitem{CE17}
D.~Cordero‑Erausquin,\\
\textit{Transport inequalities for log‑concave measures, quantitative forms, and applications},\\
Canad.~J.~Math.~69 (2017), 481--501.
\end{bibliography}
\end{document}
